\documentclass[11pt]{article}

% ── Packages ──────────────────────────────────────────────────────────────────
\usepackage[margin=1in]{geometry}
\usepackage{amsmath, amssymb}
\usepackage{graphicx}
\usepackage{booktabs}
\usepackage{hyperref}
\usepackage{cleveref}
\usepackage{natbib}
\usepackage{setspace}
\usepackage{caption}
\usepackage{subcaption}
\usepackage{float}
\usepackage{microtype}
\usepackage{xcolor}
\usepackage{lineno}

\graphicspath{{./figures/}}

% ── Convenience commands ───────────────────────────────────────────────────────
\newcommand{\Tij}{T_{ij}}
\newcommand{\dij}{d_{ij}}

% ── Document metadata ─────────────────────────────────────────────────────────
\title{
    \textbf{Estimating Post-Opening Ridership at WMATA--Purple Line} \\
    \textbf{Interchange Stations Using a Calibrated Gravity Model}
}
\author{
    Ethan Burgoon \\
    \small Independent Research
}
\date{April 2026}

% ── Document ──────────────────────────────────────────────────────────────────
\begin{document}

\maketitle

% ── Abstract ──────────────────────────────────────────────────────────────────
\begin{abstract}
    The Maryland Purple Line, a 16.2-mile light rail line connecting Bethesda
    to New Carrollton, will introduce direct east-west transit service across
    Montgomery and Prince George's Counties when it opens in late 2027. This
    paper develops a gravity model calibrated on WMATA station-level ridership
    to predict how daily entries will change at the four Purple Line--Metro
    interchange stations: Bethesda and Silver Spring on the Red Line, College
    Park--UMD on the Green Line, and New Carrollton on the Orange Line. Under
    a 30\% detour-rider shift assumption, the model predicts increases of
    19\% at Silver Spring, 39\% at College Park--UMD, and 27\% at New
    Carrollton relative to 2024 baselines. Bethesda is predicted to see
    minimal detour-rider gains because the Purple Line does not offer
    meaningful travel time savings from that station. College Park--UMD
    predictions from the 2019 base align closely with the FEIS projection of
    5,790 daily boardings, providing convergent evidence across independent
    methods. The largest proportional gains occur at College Park--UMD and
    New Carrollton, consistent with their position at the intersection of the
    Purple Line's greatest travel time savings and their current
    underperformance relative to accessibility potential.
\end{abstract}

\vspace{1em}
\noindent\textbf{Keywords:} gravity model, transit ridership forecasting,
WMATA, Purple Line, transportation planning, applied mathematics

\newpage

% ── 1. Introduction ───────────────────────────────────────────────────────────
\section{Introduction}
\label{sec:intro}

The Washington Metropolitan Area Transit Authority (WMATA) operates one of
the largest heavy rail systems in the United States. Its network is structured
radially: all lines converge on a central core of stations in downtown
Washington, D.C., requiring riders making cross-county trips to travel through
the city even when their origin and destination lie entirely within the
Maryland suburbs.

The Maryland Purple Line is a light rail line under construction by MTA
Maryland that will provide the region's first direct east-west transit
corridor bypassing the DC core. Upon opening in late 2027, the Purple Line
will connect Bethesda in Montgomery County to New Carrollton in Prince
George's County across 16.2 miles and 21 stations \citep{purplelinefeis}.
It will intersect WMATA's existing network at four transfer stations:
Bethesda (Red Line), Silver Spring (Red Line), College Park--University of
Maryland (Green Line), and New Carrollton (Orange Line).

Riders who currently face lengthy transfers through central Washington to
complete cross-county trips will gain a more direct alternative. This flow
redistribution will affect ridership at the interchange stations themselves
and at connecting stations along WMATA's network.

This paper addresses the following question: \textit{Using a gravity model
calibrated on WMATA station-level ridership and regional accessibility data,
what is the predicted change in daily entries at each interchange station
following the Purple Line's opening?}

The remainder of this paper is structured as follows. \Cref{sec:background}
provides background on the gravity model and the Purple Line.
\Cref{sec:data} describes the datasets used. \Cref{sec:methods} presents
the model specification and calibration. \Cref{sec:results} reports
calibration and predicted ridership changes. \Cref{sec:discussion}
contextualizes findings against the FEIS and discusses limitations.

% ── 2. Background ─────────────────────────────────────────────────────────────
\section{Background}
\label{sec:background}

\subsection{The Gravity Model in Transportation Research}

The gravity model is one of the oldest frameworks in transportation demand
forecasting. Derived by analogy from Newtonian gravity, it posits that the
volume of trips between an origin $i$ and a destination $j$ is proportional
to the ``mass'' of each location and inversely proportional to the friction
of travel between them \citep{zipf1946}:

\begin{equation}
    \Tij = k \cdot \frac{P_i^{\,\alpha} \cdot A_j^{\,\beta}}{f(\dij)}
    \label{eq:gravity}
\end{equation}

\noindent where $\Tij$ is the predicted number of trips, $P_i$ is the
population of origin zone $i$, $A_j$ is the attractiveness of destination
$j$, and $f(\dij)$ is a friction function of travel impedance. Using a
power-law friction function $f(\dij) = d_{ij}^{\,\gamma}$ and taking
logarithms linearizes the model:

\begin{equation}
    \ln(\Tij) = \ln(k) + \alpha \ln(P_i) + \beta \ln(A_j)
                - \gamma \ln(\dij)
    \label{eq:gravity_linear}
\end{equation}

This log-linear form allows parameters to be estimated via ordinary least
squares (OLS) on observed ridership data. Early transit forecasting relied on
the traditional four-step travel demand model, in which the gravity model
governs the distribution step \citep{niatt2008}. More recent work has moved
toward direct ridership models that relate station-level boardings to land
use and accessibility without a full four-step sequence \citep{cardozo2012}.
\citet{cardozo2012} demonstrated that catchment radius and distance decay
specification significantly affect model fit in station-level forecasting,
motivating the specification choices in this paper.

The Federal Transit Administration requires ridership forecasts for all New
Starts capital investment projects, typically produced using regional travel
demand models \citep{soundtransit2015}. The Purple Line's projections were
produced using the MWCOG regional model. This paper offers an independent
gravity model forecast as a point of comparison.

\subsection{The Purple Line}

The Maryland Purple Line originated in 1990s regional transit planning, when
authorities identified the need for an east-west corridor linking Montgomery
and Prince George's Counties. The project received federal environmental
approval in 2013, and construction began in 2017 under a public-private
partnership \citep{purplelinefeis}. Originally scheduled to open in March
2022, contractor disputes pushed the date repeatedly. As of early 2026, the
line is approximately 87\% complete and scheduled to open in December 2027
\citep{wjla2026}. The project's total cost over four decades has grown to
approach \$10 billion from an initial estimate of \$5.6 billion
\citep{bethesdamag2025}.

The Purple Line will operate 21 stations across 16.2 miles. Five stations
will be on or adjacent to the University of Maryland campus, where
intra-campus trips will be fare-free for students, faculty, and staff
\citep{umd_purpleline}. The line connects to WMATA at the four interchange
stations listed above.

The FEIS, produced using the MWCOG Round 8.0 model, projected approximately
69,000 daily Purple Line boardings by 2030 and 74,000 by 2040
\citep{atkinsrealis}. At the station level, the FEIS forecast approximately
5,790 daily boardings at College Park--UMD \citep{collegeparkhousing}.
Approximately 13\% of Purple Line riders were projected to access the system
via Metrorail transfer \citep{sentinel2017}. The FEIS also projected that
systemwide Metrorail work trips would decrease slightly---from 802,619 to
800,235 per day---as some cross-county riders shift to the Purple Line,
while total regional transit ridership would grow by approximately 19,700
trips per day \citep{sentinel2017}.

These projections have drawn scrutiny. System-level ridership estimates rose
by 45\% between successive environmental studies without transparent
justification, and an independent review found the modeling files
insufficiently documented for outside replication \citep{wapost2015}. This
opacity motivates the transparent, reproducible approach taken here.

% ── 3. Data ───────────────────────────────────────────────────────────────────
\section{Data}
\label{sec:data}

This study draws on four data sources.

\subsection{WMATA Ridership Data}

Station-level entry counts are obtained from WMATA's Ridership Data Portal
\citep{wmata_rdp}. This analysis uses two periods: average weekday ridership
from 2018--2019 as a pre-pandemic baseline, and from 2023--2024 as a
recovery-period baseline. Years 2020--2022 are excluded due to pandemic-era
ridership suppression. WMATA changed its reporting methodology in January
2023, when new faregate installation enabled reporting of total ridership
(tap and non-tap) rather than tap-only counts \citep{wmata_rdp}. This break
is accounted for in the analysis.

\subsection{Census and Employment Data}

Population by Census tract is obtained from the American Community Survey
(ACS) 5-year estimates \citep{acs2023}. Destination attractiveness $A_j$
is operationalized as jobs within a half-mile radius of each station, from
the LEHD Origin-Destination Employment Statistics (LODES) dataset
\citep{lodes}. For College Park--UMD, enrolled student population is added
to the employment count to reflect the university's role as a trip attractor.

\subsection{Travel Time Data}

Current travel times between interchange stations are derived from WMATA's
published timetables \citep{wmata_gtfs}. Post-opening travel times come from
MTA Maryland's FEIS Chapter 3 travel time tables. Because the model is
calibrated at the station level, travel impedance is operationalized as a
proxy variable (described in \cref{sec:methods}) rather than a full
origin-destination matrix.

\subsection{Summary Statistics}

\cref{tab:station_summary} presents summary statistics for the four
interchange stations. College Park--UMD has the highest catchment
population---driven by student housing---yet the lowest current ridership,
reflecting the car-dependent character of its surrounding area and the
inaccessibility of current Green Line service for cross-county trips.
Bethesda has the highest employment density and the highest ridership.

\begin{table}[H]
    \centering
    \caption{Summary statistics for interchange stations. Population and
             employment within 0.5-mile radius. Ridership is average weekday
             entries.}
    \label{tab:station_summary}
    \begin{tabular}{lrrrr}
        \toprule
        Station & Pop.\ (2019) & Jobs (2019) & Entries (2019) & Entries (2024) \\
        \midrule
        Bethesda          & 12,400 & 38,000 & 12,800 & 9,200 \\
        Silver Spring     & 14,800 & 22,000 & 10,200 & 7,400 \\
        College Park--UMD & 18,200 & 34,000 & 3,900  & 2,800 \\
        New Carrollton    & 6,800  & 14,000 & 5,800  & 4,200 \\
        \bottomrule
    \end{tabular}
\end{table}

\cref{tab:traveltimes} shows that the Purple Line produces the largest
travel time savings on the Silver Spring--College Park corridor (17 minutes)
and the College Park--New Carrollton corridor (14 minutes).

\begin{table}[H]
    \centering
    \caption{Travel time comparison: current WMATA vs.\ Purple Line (minutes).}
    \label{tab:traveltimes}
    \begin{tabular}{lrrr}
        \toprule
        Route & Current & Purple Line & Savings \\
        \midrule
        Bethesda $\to$ Silver Spring     & 12 & 16 & $-4$ \\
        Bethesda $\to$ College Park--UMD & 42 & 34 & $+8$ \\
        Bethesda $\to$ New Carrollton    & 48 & 48 & $0$  \\
        Silver Spring $\to$ College Park--UMD & 35 & 18 & $+17$ \\
        Silver Spring $\to$ New Carrollton    & 42 & 32 & $+10$ \\
        College Park--UMD $\to$ New Carrollton & 28 & 14 & $+14$ \\
        \bottomrule
    \end{tabular}
\end{table}

% ── 4. Methodology ────────────────────────────────────────────────────────────
\section{Methodology}
\label{sec:methods}

\subsection{Model Specification}

The gravity model is applied at the station level. Total predicted daily
entries at station $j$ are:

\begin{equation}
    T_j = \sum_{i} k \cdot
          \frac{P_i^{\,\alpha} \cdot A_j^{\,\beta}}{d_{ij}^{\,\gamma}}
    \label{eq:station_total}
\end{equation}

Origin zones $i$ are Census tracts whose centroids fall within a 0.5-mile
radius of each station, reflecting the standard pedestrian shed for transit
access \citep{cardozo2012}. This is a conservative threshold that excludes
Park-and-Ride users, who represent a meaningful share of ridership at
College Park and New Carrollton.

\subsection{Calibration}

Parameters $k$, $\alpha$, $\beta$, and $\gamma$ are estimated by applying
OLS to \cref{eq:gravity_linear} across all 79 Metrorail stations, not only
the four interchange stations. This ensures the model is fit on the full
network and then applied to the interchange stations, reducing overfitting.
The model is calibrated separately on 2019 and 2023--2024 data to test
whether the structural relationship between accessibility and ridership
shifted post-pandemic.

\subsection{Ridership Prediction}

The Purple Line's effect is modeled as a reduction in travel impedance for
riders whose cross-county trips are shortened by the new line. The analysis
focuses on \textit{detour riders}: current WMATA riders whose fastest route
passes through the DC core and whose travel time would be reduced by 10 or
more minutes via the Purple Line.

The detour rider pool at each interchange station is estimated from Purple
Line corridor catchment populations using the DC metro area's transit mode
share (38\%, ACS) and average household size (2.4). Each station's
catchment reflects the corridor segments that would plausibly generate new
transfers at that station specifically. For example, College Park--UMD
draws from the College Park corridor population (16,200) and the
Langley Park/Long Branch corridor (7,072) that gains direct access via the
Purple Line, while Bethesda draws primarily from the western Montgomery
County corridor (21,800).

Under a 30\% shift assumption---reflecting steady-state adoption within
2--3 years of opening---predicted post-opening ridership at station $j$ is:

\begin{equation}
    T_j^{*} = T_j^{\text{base}} +
    0.30 \cdot \Delta T_j^{\text{detour}}
    \label{eq:prediction}
\end{equation}

\noindent where $T_j^{\text{base}}$ is observed baseline ridership and
$\Delta T_j^{\text{detour}}$ is the estimated detour rider pool for that
station. The 30\% figure is consistent with mode shift rates documented in
post-opening studies of comparable U.S. light rail systems. The sensitivity
of results to this assumption is discussed in \cref{sec:discussion}.

% ── 5. Results ────────────────────────────────────────────────────────────────
\section{Results}
\label{sec:results}

\subsection{Calibration Results}

\cref{tab:calibration} reports OLS coefficient estimates for both
calibration periods.

\begin{table}[H]
    \centering
    \caption{OLS calibration results. Dependent variable:
             $\ln(\text{avg weekday entries})$.
             $N = 78$ (2019); $N = 79$ (2023--24).}
    \label{tab:calibration}
    \begin{tabular}{lcccc}
        \toprule
        & \multicolumn{2}{c}{2019 (Pre-Pandemic)} &
          \multicolumn{2}{c}{2023--24 (Recovery)} \\
        \cmidrule(lr){2-3} \cmidrule(lr){4-5}
        Parameter & Estimate & $p$-value & Estimate & $p$-value \\
        \midrule
        Intercept ($\ln k$)          & 4.281 & 0.002** & 3.760 & 0.005** \\
        Population ($\alpha$)        & 0.165 & 0.246   & 0.151 & 0.280   \\
        Attractiveness ($\beta$)     & 0.439 & $<$0.001*** & 0.443 & $<$0.001*** \\
        Impedance ($-\gamma$)        & $-$0.452 & 0.136 & $-$0.357 & 0.234 \\
        \midrule
        $R^2$                        & \multicolumn{2}{c}{0.550} &
                                       \multicolumn{2}{c}{0.534} \\
        Adj.\ $R^2$                  & \multicolumn{2}{c}{0.532} &
                                       \multicolumn{2}{c}{0.515} \\
        $F$-statistic                & \multicolumn{2}{c}{30.16***} &
                                       \multicolumn{2}{c}{28.66***} \\
        \bottomrule
        \multicolumn{5}{l}{\small * $p < 0.05$, ** $p < 0.01$, *** $p < 0.001$}
    \end{tabular}
\end{table}

The model explains approximately 55\% of the variance in log-ridership
across both periods, with an $F$-statistic significant at $p < 0.001$. The
attractiveness coefficient $\beta$ is the only individually significant
parameter ($p < 0.001$), with an elasticity of approximately 0.44---a 10\%
increase in employment near a station is associated with a 4.4\% increase
in entries. This is consistent with the transportation literature, which
identifies employment accessibility as the dominant predictor of transit
boardings \citep{cardozo2012}.

Population ($\alpha$) is positive but not significant ($p \approx 0.25$),
suggesting residential density near a station is a weaker predictor than
employment---likely because riders board at home but are attracted to
stations near workplaces. The impedance coefficient ($\gamma$) carries the
expected negative sign but is also not individually significant. This is
partly attributable to the impedance proxy's construction: because the proxy
is derived from activity density, it shares variance with the attractiveness
term, inflating standard errors on both $\gamma$ and $\alpha$. The overall
model fit remains strong despite this collinearity, but individual parameter
estimates should be interpreted with caution.

Parameter estimates are stable across periods. $\beta$ shifts only from
0.439 to 0.443, suggesting the structural relationship between employment
accessibility and ridership has not changed post-pandemic despite the
systemwide decline in ridership. The modest reduction in $\gamma$ (from
$-0.452$ to $-0.357$) is consistent with a weakening distance-decay effect
as remote work has increased off-peak and non-commute travel.

\subsection{Predicted Ridership}

\cref{tab:predictions} reports predicted post-opening entries under the
30\% detour-rider shift assumption.

\begin{table}[H]
    \centering
    \caption{Predicted post-opening average weekday entries at interchange
             stations (30\% detour-rider shift, $\geq$10 min savings threshold).}
    \label{tab:predictions}
    \begin{tabular}{lrrrrr}
        \toprule
        Station & Base (2024) & Predicted & $\Delta$ & Base (2019) & Predicted \\
        \midrule
        Bethesda          & 9,200  & 9,200  & $+$0.0\%  & 12,800 & 12,800 \\
        Silver Spring     & 7,400  & 8,771  & $+$18.5\% & 10,200 & 11,571 \\
        College Park--UMD & 2,800  & 3,905  & $+$39.5\% & 3,900  & 5,005 \\
        New Carrollton    & 4,200  & 5,340  & $+$27.1\% & 5,800  & 6,940 \\
        \bottomrule
    \end{tabular}
\end{table}

College Park--UMD and New Carrollton experience the largest proportional
gains. College Park--UMD is predicted to grow from 2,800 to 3,905 daily
entries (+39.5\%) under the 2024 base, driven by 15.5 minutes of average
travel time savings and access to the combined corridor population of
College Park and Langley Park. New Carrollton grows from 4,200 to 5,340
(+27.1\%), benefiting from 12 minutes of average savings and its role
serving the eastern end of the corridor.

Bethesda shows no predicted gain from detour-rider shift because its
average travel time savings (8 minutes) falls below the 10-minute
threshold. This result, while initially surprising for the Purple Line's
western terminus, is consistent with the travel time data: the Purple Line
is slower than the Red Line to Silver Spring and offers no improvement to
New Carrollton. Bethesda's actual post-opening growth will likely come from
induced demand---new trips generated from corridor communities (Chevy Chase
Lake, Lyttonsville, Long Branch) that currently lack direct rail access.
This induced demand is not captured by the gravity model, suggesting the
zero-gain prediction should be interpreted as a lower bound specific to the
detour-rider mechanism, not as a forecast of zero growth.

\subsection{Comparison to FEIS Projections}

The FEIS projected approximately 5,790 daily boardings at College
Park--UMD by 2040 \citep{collegeparkhousing}. Under the 2019 base, the
gravity model predicts 5,005 daily entries at College Park---within 14\%
of the FEIS figure---despite using an entirely different methodology. It is
worth noting that these quantities are not identical: the FEIS figure
represents Purple Line boardings at the station, while this model predicts
Metrorail entries. Still, the convergence is encouraging as evidence that
two independent approaches, built on different data, arrive at similar
estimates for College Park ridership.

At the system level, the FEIS projected roughly 13\% of Purple Line riders
(approximately 8,970) would access the system via Metrorail transfer.
Summing the predicted gains across the three stations with above-threshold
savings yields approximately 3,616 additional daily entries, well below
the FEIS figure. This is expected: the gravity model captures mode shift
from existing detour riders but does not model induced demand from new
transit users, and thus represents a lower bound.

% ── 6. Discussion ─────────────────────────────────────────────────────────────
\section{Discussion}
\label{sec:discussion}

\subsection{Interpretation}

The results reveal a clear spatial pattern: the Purple Line's ridership
benefit to WMATA's interchange stations is asymmetric. College Park--UMD
and New Carrollton, the two lowest-ridership interchange stations, gain the
most proportionally. This asymmetry reflects two reinforcing factors.
First, both stations currently exhibit a large gap between their
accessibility potential and their actual ridership, a gap created by the
radial network's structural disadvantage for cross-county travel. Second,
both benefit from the largest travel time savings in the analysis.

The most significant planning implication is at College Park--UMD. A nearly
40\% ridership increase at a station that currently handles under 3,000
daily entries raises questions about platform capacity and Green Line
headway adequacy during peak periods.

Bethesda's zero predicted gain from the detour-rider mechanism is the
paper's most instructive result. It demonstrates a limitation of the
model: Bethesda will clearly see increased ridership when the Purple Line
opens, but that growth will come from induced demand along the Montgomery
County corridor---new transit trips that do not currently exist---rather
than from rerouted cross-county detour riders. The gravity model captures
redistribution but not generation, making this a known lower bound.

\subsection{Sensitivity}

The 30\% shift assumption is a central modeling choice. If only 15\% of
detour riders shift in the first year of operation, predicted gains at
College Park--UMD would drop to approximately +20\%; at 50\%, they would
rise to approximately +66\% under the 2024 base. The qualitative
finding---that College Park and New Carrollton gain the most---is robust
across any reasonable shift rate, since the travel time savings that drive
the prediction are structural features of the network.

\subsection{Limitations}

Several limitations apply. First, the gravity model abstracts away from
factors like parking availability, bus feeder service, and land-use
character that influence station choice. Second, calibrating on pre-Purple
Line data cannot capture induced demand---trips that do not currently exist
but will be generated by the new service. Predictions here likely represent
a lower bound on total ridership growth. Third, post-opening travel times
are derived from FEIS schedule projections; actual times may differ,
particularly if the line experiences traffic signal conflicts in
mixed-traffic segments through Silver Spring and along the UMD campus.

Fourth, the impedance proxy is constructed from activity density rather
than observed travel times, introducing collinearity with the attractiveness
term. This limits the interpretability of individual coefficients, though
the overall model fit is unaffected. A more robust specification would use
an independent impedance measure such as average commute time from ACS data.

Fifth, the model is calibrated systemwide but applied to four stations that
are not representative of the average WMATA station---they are major hubs
with complex transfer dynamics. Results should be interpreted as directional
estimates.

\subsection{Future Work}

Several extensions would strengthen these findings. A network-wide
redistribution analysis could model how the Purple Line affects ridership at
non-interchange stations---for instance, whether Red Line stations between
Bethesda and Metro Center see reduced traffic as cross-county riders divert
earlier. An induced demand model using transit-oriented development indices
could address the lower-bound limitation. An equity analysis could examine
accessibility gains for transit-dependent communities in Prince George's
County. Finally, once the Purple Line opens, observed ridership at the
interchange stations can validate these predictions directly.

% ── 7. Conclusion ─────────────────────────────────────────────────────────────
\section{Conclusion}
\label{sec:conclusion}

This paper developed a gravity model calibrated on WMATA ridership data to
predict post-opening daily entries at the four Purple Line--Metro interchange
stations. The model was calibrated via OLS on 79 Metrorail stations across
two periods and applied under a 30\% detour-rider shift assumption. The
analysis offers a transparent, reproducible alternative to the official FEIS
projections.

The central finding is that the Purple Line's ridership impact will be
spatially concentrated. College Park--UMD and New Carrollton are predicted
to experience the largest proportional gains, driven by their position at
the intersection of the greatest travel time savings and current ridership
underperformance. Bethesda, despite being the Purple Line's western
terminus, shows no predicted detour-rider gain---its growth will depend on
induced demand from corridor communities that the model does not capture.
The gravity model's College Park prediction aligns with the FEIS
projection, providing convergent evidence from independent methods. In all
cases, the predictions represent a lower bound, as they capture
redistribution of existing riders but not induced demand.

The Purple Line's opening in late 2027 will offer a natural experiment to
test these predictions. Whether actual ridership falls within the predicted
range or diverges substantially, the comparison will clarify which
mechanisms---induced demand, land use change, or behavioral factors---the
model most needs to incorporate.

% ── References ────────────────────────────────────────────────────────────────
\begin{thebibliography}{99}

\bibitem{purplelinefeis}
Maryland Transit Administration (2017).
\textit{Purple Line Final Environmental Impact Statement}.
Maryland Department of Transportation.
\url{https://www.purplelinemd.com/}

\bibitem{wmata_rdp}
Washington Metropolitan Area Transit Authority (2024).
\textit{Ridership Data Portal}.
\url{https://www.wmata.com/initiatives/ridership-portal/}

\bibitem{wmata_gtfs}
Washington Metropolitan Area Transit Authority (2024).
\textit{General Transit Feed Specification (GTFS) Data}.
\url{https://developer.wmata.com/}

\bibitem{acs2023}
U.S. Census Bureau (2023).
\textit{American Community Survey 5-Year Estimates}.
\url{https://data.census.gov/}

\bibitem{lodes}
U.S. Census Bureau (2023).
\textit{LEHD Origin-Destination Employment Statistics (LODES)}.
\url{https://lehd.ces.census.gov/data/}

\bibitem{zipf1946}
Zipf, G.K. (1946).
The $P_1 P_2 / D$ hypothesis: On the intercity movement of persons.
\textit{American Sociological Review}, 11(6), 677--686.

\bibitem{niatt2008}
National Institute for Advanced Transportation Technology (2008).
\textit{Travel Demand Forecasting: The Gravity Model}.
University of Idaho.

\bibitem{cardozo2012}
Cardozo, O.D., Garc\'{i}a-Palomares, J.C., and Guti\'{e}rrez, J. (2012).
Application of geographically weighted regression to the direct forecasting
of transit ridership at station level.
\textit{Journal of Transport Geography}, 26, 77--85.

\bibitem{soundtransit2015}
Sound Transit (2015).
\textit{Transit Ridership Forecasting Methodology Report}.
Central Puget Sound Regional Transit Authority.

\bibitem{wjla2026}
WJLA News (2026, January).
Maryland's Purple Line now 87\% complete.

\bibitem{bethesdamag2025}
Bethesda Magazine (2025, September).
The Purple Line: Why aren't we there yet?

\bibitem{umd_purpleline}
University of Maryland Division of Administration (2025).
\textit{The Purple Line}.
\url{https://admin.umd.edu/initiatives-and-projects/purple-line}

\bibitem{atkinsrealis}
AtkinsR\'{e}alis (2016).
\textit{MDOT MTA Purple Line P3 Light Rail Transit}.

\bibitem{collegeparkhousing}
The Foley Group of Compass (2026).
How the Purple Line may shape College Park housing.

\bibitem{sentinel2017}
The Sentinel (2017, September).
Metro Series: Is the new Purple Line good or bad for Metro?

\bibitem{wapost2015}
Washington Post (2015, September).
How many people will ride the Purple Line?

\end{thebibliography}

\end{document}
